\documentclass[11pt]{article}
% ===========================================================================
%  Shared preamble for all MPM2D worksheets — input right after \documentclass.
%  (Files beginning with "_" are skipped by mpm2d-worksheet-publish.mjs.)
% ===========================================================================
\usepackage[letterpaper,top=2.2cm,bottom=1.9cm,left=1.6cm,right=1.6cm,headheight=28pt,footskip=24pt]{geometry}
\usepackage{amsmath,amssymb}
\usepackage{cancel}
\usepackage[T1]{fontenc}
\usepackage[utf8]{inputenc}
\usepackage{lmodern}
\usepackage{xcolor}
\usepackage[most]{tcolorbox}
\usepackage{enumitem}
\usepackage{multicol}
\usepackage{fancyhdr}
\usepackage{tikz}
\usetikzlibrary{calc}
\usepackage{pgfplots}
\pgfplotsset{compat=1.18}

\definecolor{exblue}{HTML}{4A90E2}
\definecolor{exbg}{HTML}{E6F3FF}
\definecolor{qorange}{HTML}{E69138}
\definecolor{qbg}{HTML}{FFF7CC}
\definecolor{keygray}{HTML}{475569}
\definecolor{keybg}{HTML}{F1F5F9}
\definecolor{iagreen}{HTML}{14653B}
\definecolor{titleblue}{HTML}{1D4ED8}
% Rotating palette so each worked example gets its own colour.
\definecolor{ex0}{HTML}{2563A0}\definecolor{ex1}{HTML}{3B7D3B}\definecolor{ex2}{HTML}{8A5A00}
\definecolor{ex3}{HTML}{A3327A}\definecolor{ex4}{HTML}{6D28A3}\definecolor{ex5}{HTML}{B08900}
\definecolor{ex6}{HTML}{0E7490}\definecolor{ex7}{HTML}{9A3412}\definecolor{ex8}{HTML}{15803D}

\pagestyle{fancy}
\fancyhf{}
\fancyhead[L]{\footnotesize NAME: \underline{\hspace{3.4cm}}}
\fancyhead[C]{\textbf{Integration Academy}\\[-2pt]{\scriptsize Math Simplified \textperiodcentered{} Success Amplified}}
\fancyhead[R]{\footnotesize MPM2D}
\fancyfoot[L]{\scriptsize dr.merie@IntegrationAcademy.ca}
\fancyfoot[C]{\scriptsize\thepage}
\fancyfoot[R]{\scriptsize IntegrationAcademy.ca}
\renewcommand{\headrulewidth}{1.2pt}
\renewcommand{\footrulewidth}{0.4pt}
\setlength{\parindent}{0pt}
\setlength{\parskip}{4pt}

% Examples are NOT breakable, so each example + its solution stays on one page.
\newtcolorbox{example}[2]{colframe=#1,colback=#1!7,coltitle=white,colbacktitle=#1,fonttitle=\bfseries,title={#2},arc=3pt,boxrule=1pt}
\newtcolorbox{qbox}[1]{colframe=qorange,colback=qbg,coltitle=white,colbacktitle=qorange,fonttitle=\bfseries,title={#1},arc=3pt,breakable,boxrule=1pt}
\newtcolorbox{keybox}{colframe=keygray,colback=keybg,coltitle=white,colbacktitle=keygray,fonttitle=\bfseries,title={$\checkmark$\ Answer Key},arc=3pt,breakable,boxrule=1pt}
\newtcolorbox{recapbox}{colframe=keygray,colback=keybg,coltitle=white,colbacktitle=keygray,fonttitle=\bfseries,title={Question Recap},arc=3pt,breakable,boxrule=1pt}
\newtcolorbox{ideabox}{colframe=titleblue,colback=titleblue!6,coltitle=white,colbacktitle=titleblue,fonttitle=\bfseries,title={The Big Ideas},arc=3pt,breakable,boxrule=1.2pt}
\newtcolorbox{learnbox}{colframe=iagreen,colback=iagreen!5,coltitle=white,colbacktitle=iagreen,fonttitle=\bfseries,title={Concept Spotlight},arc=3pt,breakable,boxrule=1.2pt}
\newcommand{\lh}[1]{\par\smallskip{\bfseries\color{iagreen}#1}\par\nobreak\smallskip}

\newcommand{\soln}{\par\smallskip\textit{\textbf{Solution.}}\ }
\newcommand{\workspace}[1]{\par\vspace{3pt}\fcolorbox{black!30}{white}{\begin{minipage}[t][#1][t]{\dimexpr\linewidth-2\fboxsep\relax}\ \end{minipage}}\par}
\newcommand{\sech}[1]{\par\vspace{8pt}{\Large\bfseries\color{iagreen}#1}\par\nobreak\vspace{2pt}{\color{black!20}\hrule height1pt}\vspace{6pt}}

% A small coordinate plot: \plot{xmin}{xmax}{ymin}{ymax}{body}
\newcommand{\plot}[5]{\begin{center}\begin{tikzpicture}
\begin{axis}[width=8.4cm,height=6cm,axis lines=middle,xlabel={\scriptsize$x$},ylabel={\scriptsize$y$},
grid=both,grid style={gray!16},xmin=#1,xmax=#2,ymin=#3,ymax=#4,
xtick distance=2,ytick distance=2,tick label style={font=\tiny},axis line style={-}]
#5
\end{axis}\end{tikzpicture}\end{center}}

% Right triangle (angle theta at bottom-left, right angle at bottom-right):
%   \rtri{drawBase}{drawHeight}{oppLabel}{adjLabel}{hypLabel}
\newcommand{\rtri}[5]{\begin{center}\begin{tikzpicture}[scale=0.85]
\coordinate (A) at (0,0);\coordinate (B) at (#1,0);\coordinate (C) at (#1,#2);
\draw[thick,exblue] (A)--(B)--(C)--cycle;
\draw[black] ($(B)+(-0.32,0)$)--($(B)+(-0.32,0.32)$)--($(B)+(0,0.32)$);
\node[below] at ($(A)!0.5!(B)$){#4};
\node[right] at ($(B)!0.5!(C)$){#3};
\node[above left] at ($(A)!0.5!(C)$){#5};
\node at (0.7,0.22){$\theta$};
\end{tikzpicture}\end{center}}

% General (oblique) triangle with vertex labels and opposite-side labels:
%   \gtri{Alabel}{Blabel}{Clabel}{aSide}{bSide}{cSide}
\newcommand{\gtri}[6]{\begin{center}\begin{tikzpicture}[scale=0.85]
\coordinate (A) at (0,0);\coordinate (B) at (5,0);\coordinate (C) at (1.6,3);
\draw[thick,exblue] (A)--(B)--(C)--cycle;
\node[below left] at (A){#1};\node[below right] at (B){#2};\node[above] at (C){#3};
\node[below] at ($(A)!0.5!(B)$){#6};
\node[right] at ($(B)!0.5!(C)$){#4};
\node[left] at ($(A)!0.5!(C)$){#5};
\end{tikzpicture}\end{center}}

\fancyhead[R]{\footnotesize MPM2D \textperiodcentered{} 1.4}
\begin{document}

\begin{center}
{\LARGE\bfseries\color{titleblue}Worksheet 1.4 --- Solving by Elimination}\\[2pt]
{\small Grade 10 Academic (MPM2D) \textperiodcentered{} Unit 1: Linear Systems}
\end{center}

\begin{ideabox}
\textbf{Elimination} adds or subtracts the equations so that one variable cancels.
\begin{itemize}[leftmargin=*,itemsep=2pt]
  \item Line the equations up as $Ax+By=C$.
  \item If a variable's coefficients are \emph{opposites}, \textbf{add}; if they are \emph{equal}, \textbf{subtract}.
  \item If neither matches, \textbf{multiply} one (or both) equations to create a matching pair first.
  \item Solve for the remaining variable, then back-substitute.
\end{itemize}
\end{ideabox}

\begin{learnbox}
\lh{The idea}\textbf{Elimination} adds or subtracts the two equations to cancel one variable. If needed, multiply an equation first so a pair of coefficients line up.
\lh{The steps}Stack the equations, scale one or both so a variable's coefficients are equal or opposite, add or subtract to eliminate it, solve, then back-substitute.
\lh{Choosing a method}Elimination is often fastest when no variable is already isolated, especially when the coefficients share a simple multiple.
\end{learnbox}

\sech{Worked Examples}

\begin{example}{ex0}{Example 1 --- Add to eliminate}
Solve:\[\begin{array}{rll} x+y & = 10 & \text{\textcircled{\scriptsize 1}}\\ x-y & = 4 & \text{\textcircled{\scriptsize 2}} \end{array}\]\soln
\textbf{Step 1:} Add the equations to eliminate $y$:\[\begin{array}{rl} & x+\cancel{y}=10\\ +\,( & x-\cancel{y}=4\,)\\ \hline & (x+x)+(y-y)=10+4\\ & 2x=14\\ & \frac{2x}{2}=\frac{14}{2}\\ & x=7 \end{array}\]
\textbf{Step 2:} Substitute $x=7$ into \textcircled{\scriptsize 1}:\[\begin{array}{l} 7+y=10\\ y=10-7\\ y=3 \end{array}\]
\textbf{Conclusion:} $(x,y)=(7,3)$.
Graphing $y=10-x$ and $y=x-4$ shows the lines crossing at $(7,3)$.
\plot{-1}{12}{-6}{12}{\addplot[exblue,very thick,domain=-1:12,samples=2]{10-x};\addplot[qorange,very thick,domain=-1:12,samples=2]{x-4};\addplot[mark=*,only marks,ex3]coordinates{(7,3)};}
\end{example}

\begin{example}{ex1}{Example 2 --- Add (opposite $y$)}
Solve:\[\begin{array}{rll} 2x+y & = 7 & \text{\textcircled{\scriptsize 1}}\\ x-y & = 2 & \text{\textcircled{\scriptsize 2}} \end{array}\]\soln
\textbf{Step 1:} Add the equations to eliminate $y$:\[\begin{array}{rl} & 2x+\cancel{y}=7\\ +\,( & x-\cancel{y}=2\,)\\ \hline & (2x+x)+(y-y)=7+2\\ & 3x=9\\ & \frac{3x}{3}=\frac{9}{3}\\ & x=3 \end{array}\]
\textbf{Step 2:} Substitute $x=3$ into \textcircled{\scriptsize 2}:\[\begin{array}{l} 3-y=2\\ -y=2-3\\ -y=-1\\ \frac{-y}{-1}=\frac{-1}{-1}\\ y=1 \end{array}\]
\textbf{Conclusion:} $(x,y)=(3,1)$.
Graphing $y=7-2x$ and $y=x-2$ confirms the crossing point.
\plot{-1}{6}{-6}{10}{\addplot[exblue,very thick,domain=-1:6,samples=2]{7-2*x};\addplot[qorange,very thick,domain=-1:6,samples=2]{x-2};\addplot[mark=*,only marks,ex3]coordinates{(3,1)};}
\end{example}

\begin{example}{ex2}{Example 3 --- Opposite coefficients}
Solve:\[\begin{array}{rll} 3x+2y & = 12 & \text{\textcircled{\scriptsize 1}}\\ 3x-2y & = 0 & \text{\textcircled{\scriptsize 2}} \end{array}\]\soln
\textbf{Step 1:} Add the equations to eliminate $y$ (the $y$-coefficients are opposites):\[\begin{array}{rl} & 3x+\cancel{2y}=12\\ +\,( & 3x-\cancel{2y}=0\,)\\ \hline & (3x+3x)+(2y-2y)=12+0\\ & 6x=12\\ & \frac{6x}{6}=\frac{12}{6}\\ & x=2 \end{array}\]
\textbf{Step 2:} Substitute $x=2$ into \textcircled{\scriptsize 1}:\[\begin{array}{l} 3(2)+2y=12\\ 6+2y=12\\ 2y=12-6\\ 2y=6\\ \frac{2y}{2}=\frac{6}{2}\\ y=3 \end{array}\]
\textbf{Conclusion:} $(x,y)=(2,3)$.
Graphing $y=6-1.5x$ and $y=1.5x$ shows the lines meeting at $(2,3)$.
\plot{-1}{5}{-2}{8}{\addplot[exblue,very thick,domain=-1:5,samples=2]{6-1.5*x};\addplot[qorange,very thick,domain=-1:5,samples=2]{1.5*x};\addplot[mark=*,only marks,ex3]coordinates{(2,3)};}
\end{example}

\begin{example}{ex3}{Example 4 --- Multiply one equation}
Solve:\[\begin{array}{rll} 2x+y & = 5 & \text{\textcircled{\scriptsize 1}}\\ x+3y & = 10 & \text{\textcircled{\scriptsize 2}} \end{array}\]\soln
\textbf{Step 1:} Multiply \textcircled{\scriptsize 1} by $3$, then subtract \textcircled{\scriptsize 2}:\[\begin{array}{rl} & 6x+\cancel{3y}=15\\ -\,( & x+\cancel{3y}=10\,)\\ \hline & (6x-x)+(3y-3y)=15-10\\ & 5x=5\\ & \frac{5x}{5}=\frac{5}{5}\\ & x=1 \end{array}\]
\textbf{Step 2:} Substitute $x=1$ into \textcircled{\scriptsize 1}:\[\begin{array}{l} 2(1)+y=5\\ 2+y=5\\ y=5-2\\ y=3 \end{array}\]
\textbf{Conclusion:} $(x,y)=(1,3)$.
Graphing $y=5-2x$ and $y=\tfrac{10-x}3$ confirms the crossing point.
\plot{-1}{4}{-4}{8}{\addplot[exblue,very thick,domain=-1:4,samples=2]{5-2*x};\addplot[qorange,very thick,domain=-1:4,samples=2]{(10-x)/3};\addplot[mark=*,only marks,ex3]coordinates{(1,3)};}
\end{example}

\begin{example}{ex4}{Example 5 --- Multiply both equations}
Solve:\[\begin{array}{rll} 2x+3y & = 7 & \text{\textcircled{\scriptsize 1}}\\ 3x+2y & = 8 & \text{\textcircled{\scriptsize 2}} \end{array}\]\soln
\textbf{Step 1:} To eliminate $x$, multiply \textcircled{\scriptsize 1} by $3$ and \textcircled{\scriptsize 2} by $2$, then subtract:\[\begin{array}{rl} & \cancel{6x}+9y=21\\ -\,( & \cancel{6x}+4y=16\,)\\ \hline & (9y-4y)=(21-16)\\ & 5y=5\\ & \frac{5y}{5}=\frac{5}{5}\\ & y=1 \end{array}\]
\textbf{Step 2:} Substitute $y=1$ into \textcircled{\scriptsize 1}:\[\begin{array}{l} 2x+3(1)=7\\ 2x+3=7\\ 2x=7-3\\ 2x=4\\ \frac{2x}{2}=\frac{4}{2}\\ x=2 \end{array}\]
\textbf{Conclusion:} $(x,y)=(2,1)$.
Graphing $y=\tfrac{7-2x}3$ and $y=\tfrac{8-3x}2$ shows the lines crossing at $(2,1)$.
\plot{-1}{5}{-4}{6}{\addplot[exblue,very thick,domain=-1:5,samples=2]{(7-2*x)/3};\addplot[qorange,very thick,domain=-1:5,samples=2]{(8-3*x)/2};\addplot[mark=*,only marks,ex3]coordinates{(2,1)};}
\end{example}

\begin{example}{ex5}{Example 6 --- Infinitely many}
Solve:\[\begin{array}{rll} 4x-2y & = 2 & \text{\textcircled{\scriptsize 1}}\\ 2x-y & = 1 & \text{\textcircled{\scriptsize 2}} \end{array}\]\soln
\textbf{Step 1:} Multiply \textcircled{\scriptsize 2} by $2$, then subtract it from \textcircled{\scriptsize 1}:\[\begin{array}{rl} & \cancel{4x}-\cancel{2y}=2\\ -\,( & \cancel{4x}-\cancel{2y}=2\,)\\ \hline & 0=2-2\\ & 0=0 \end{array}\]

\textbf{Conclusion:} $0=0$ is always true, so there are infinitely many solutions (both equations describe the same line).
Graphing $4x-2y=2$ as $y=2x-1$ shows both equations are the exact same line.
\plot{-2}{4}{-6}{8}{\addplot[exblue,very thick,domain=-2:4,samples=2]{2*x-1};}
\end{example}

\begin{example}{ex6}{Example 7 --- No solution}
Solve:\[\begin{array}{rll} x+y & = 3 & \text{\textcircled{\scriptsize 1}}\\ 2x+2y & = 8 & \text{\textcircled{\scriptsize 2}} \end{array}\]\soln
\textbf{Step 1:} Multiply \textcircled{\scriptsize 1} by $2$, then subtract \textcircled{\scriptsize 2}:\[\begin{array}{rl} & \cancel{2x}+\cancel{2y}=6\\ -\,( & \cancel{2x}+\cancel{2y}=8\,)\\ \hline & 0=6-8\\ & 0=-2 \end{array}\]

\textbf{Conclusion:} $0=-2$ is false, so there is no solution (the lines are parallel).
Graphing $y=3-x$ and $y=4-x$ shows why: both have slope $-1$ but different $y$-intercepts, so the lines are parallel and never cross.
\plot{-2}{5}{-3}{8}{\addplot[exblue,very thick,domain=-2:5,samples=2]{3-x};\addplot[qorange,very thick,domain=-2:5,samples=2]{4-x};}
\end{example}

\begin{example}{ex7}{Example 8 --- Word problem}
Two adults and three children pay \$31; one adult and two children pay \$18. Find each price.\soln Let $a$ be the adult price and $c$ the child price, so the system is\[\begin{array}{rll} 2a+3c & = 31 & \text{\textcircled{\scriptsize 1}}\\ a+2c & = 18 & \text{\textcircled{\scriptsize 2}} \end{array}\]
\textbf{Step 1:} Multiply \textcircled{\scriptsize 2} by $2$, then subtract \textcircled{\scriptsize 1}:\[\begin{array}{rl} & \cancel{2a}+4c=36\\ -\,( & \cancel{2a}+3c=31\,)\\ \hline & (4c-3c)=(36-31)\\ & c=5 \end{array}\]
\textbf{Step 2:} Substitute $c=5$ into \textcircled{\scriptsize 2}:\[\begin{array}{l} a+2(5)=18\\ a+10=18\\ a=18-10\\ a=8 \end{array}\]
\textbf{Conclusion:} $(a,c)=(8,5)$: an adult ticket costs \$8 and a child ticket costs \$5.
Graphing $c=\tfrac{31-2a}3$ and $c=\tfrac{18-a}2$ (adult price $a$ on the horizontal axis, child price $c$ on the vertical) shows the lines crossing at $(8,5)$.
\plot{0}{15}{-1}{11}{\addplot[exblue,very thick,domain=0:15,samples=2]{(31-2*x)/3};\addplot[qorange,very thick,domain=0:15,samples=2]{(18-x)/2};\addplot[mark=*,only marks,ex3]coordinates{(8,5)};}
\end{example}

\begin{example}{ex8}{Example 9 --- Multiply both (harder)}
Solve:\[\begin{array}{rll} 3x+2y & = 4 & \text{\textcircled{\scriptsize 1}}\\ 2x+5y & = -1 & \text{\textcircled{\scriptsize 2}} \end{array}\]\soln
\textbf{Step 1:} To eliminate $y$, multiply \textcircled{\scriptsize 1} by $5$ and \textcircled{\scriptsize 2} by $2$, then subtract:\[\begin{array}{rl} & 15x+\cancel{10y}=20\\ -\,( & 4x+\cancel{10y}=-2\,)\\ \hline & (15x-4x)=(20-(-2))\\ & 11x=22\\ & \frac{11x}{11}=\frac{22}{11}\\ & x=2 \end{array}\]
\textbf{Step 2:} Substitute $x=2$ into \textcircled{\scriptsize 1}:\[\begin{array}{l} 3(2)+2y=4\\ 6+2y=4\\ 2y=4-6\\ 2y=-2\\ \frac{2y}{2}=\frac{-2}{2}\\ y=-1 \end{array}\]
\textbf{Conclusion:} $(x,y)=(2,-1)$.
Graphing $y=\tfrac{4-3x}2$ and $y=\tfrac{-1-2x}5$ confirms the crossing point.
\plot{-2}{5}{-6}{6}{\addplot[exblue,very thick,domain=-2:5,samples=2]{(4-3*x)/2};\addplot[qorange,very thick,domain=-2:5,samples=2]{(-1-2*x)/5};\addplot[mark=*,only marks,ex3]coordinates{(2,-1)};}
\end{example}

\clearpage
\sech{Your Turn --- Show Your Work}
For each question, show all steps clearly in the space provided.

\begin{qbox}{Question 1}
Solve by elimination: $x+y=8$ and $x-y=2$.\workspace{2.4cm}
\end{qbox}

\begin{qbox}{Question 2}
Solve by elimination: $2x+y=9$ and $x-y=0$.\workspace{2.4cm}
\end{qbox}

\begin{qbox}{Question 3}
Solve by elimination: $x+2y=7$ and $x-2y=1$.\workspace{2.4cm}
\end{qbox}

\begin{qbox}{Question 4}
Solve by elimination: $3x+y=10$ and $x+y=4$.\workspace{2.4cm}
\end{qbox}

\begin{qbox}{Question 5}
Solve by elimination: $2x+3y=12$ and $2x-y=4$.\workspace{2.6cm}
\end{qbox}

\begin{qbox}{Question 6}
Solve by elimination: $x+3y=9$ and $2x+y=8$.\workspace{2.6cm}
\end{qbox}

\begin{qbox}{Question 7}
Solve by elimination: $3x+2y=7$ and $5x+2y=11$.\workspace{2.4cm}
\end{qbox}

\begin{qbox}{Question 8}
Solve by elimination: $2x+5y=1$ and $3x-5y=14$.\workspace{2.4cm}
\end{qbox}

\begin{qbox}{Question 9}
Solve: $2x-y=3$ and $4x-2y=6$. How many solutions?\workspace{2.4cm}
\end{qbox}

\begin{qbox}{Question 10}
Solve: $x+y=5$ and $2x+2y=12$. How many solutions?\workspace{2.4cm}
\end{qbox}

\begin{qbox}{Question 11}
Three pens and two erasers cost \$13; one pen and two erasers cost \$7. Find each price.\workspace{2.8cm}
\end{qbox}

\begin{qbox}{Question 12}
Solve by elimination: $4x+3y=10$ and $2x+y=4$.\workspace{2.6cm}
\end{qbox}

\begin{qbox}{Question 13 --- Challenge}
Solve by elimination: $3x+4y=5$ and $2x+3y=4$.\workspace{3.0cm}
\end{qbox}

\clearpage
\sech{Question Recap \& Answer Key}

\begin{recapbox}
\begin{multicols}{2}
\small
\begin{enumerate}[leftmargin=*,itemsep=3pt]
  \item Solve by elimination: $x+y=8$ and $x-y=2$.
  \item Solve by elimination: $2x+y=9$ and $x-y=0$.
  \item Solve by elimination: $x+2y=7$ and $x-2y=1$.
  \item Solve by elimination: $3x+y=10$ and $x+y=4$.
  \item Solve by elimination: $2x+3y=12$ and $2x-y=4$.
  \item Solve by elimination: $x+3y=9$ and $2x+y=8$.
  \item Solve by elimination: $3x+2y=7$ and $5x+2y=11$.
  \item Solve by elimination: $2x+5y=1$ and $3x-5y=14$.
  \item Solve: $2x-y=3$ and $4x-2y=6$. How many solutions?
  \item Solve: $x+y=5$ and $2x+2y=12$. How many solutions?
  \item Three pens and two erasers cost \$13; one pen and two erasers cost \$7. Find each price.
  \item Solve by elimination: $4x+3y=10$ and $2x+y=4$.
  \item Solve by elimination: $3x+4y=5$ and $2x+3y=4$.
\end{enumerate}
\end{multicols}
\end{recapbox}

\begin{keybox}
\begin{multicols}{2}
\small
\begin{enumerate}[leftmargin=*,itemsep=3pt]
  \item Add: $2x=10\Rightarrow x=5$, $y=3$: $(5,3)$.
  \item Add: $3x=9\Rightarrow x=3$, $y=3$: $(3,3)$.
  \item Add: $2x=8\Rightarrow x=4$, $y=1.5$: $(4,1.5)$.
  \item Subtract: $2x=6\Rightarrow x=3$, $y=1$: $(3,1)$.
  \item Subtract: $4y=8\Rightarrow y=2$, $x=3$: $(3,2)$.
  \item $\times2$ first, subtract: $5y=10\Rightarrow y=2$, $x=3$: $(3,2)$.
  \item Subtract: $2x=4\Rightarrow x=2$, $y=0.5$: $(2,0.5)$.
  \item Add: $5x=15\Rightarrow x=3$, $y=-1$: $(3,-1)$.
  \item Infinitely many (second $=2\times$ first).
  \item No solution ($2x+2y$ can't equal both $10$ and $12$).
  \item Subtract: $2\text{ pens}=6\Rightarrow$ pen \$3, eraser \$2.
  \item $\times2$ second, subtract: $y=2$, $x=1$: $(1,2)$.
  \item $\times3,\times4$: $9x+12y=15$, $8x+12y=16$; subtract $x=-1$, $y=2$: $(-1,2)$.
\end{enumerate}
\end{multicols}
\end{keybox}

\end{document}
